\documentclass[a4paper,13pt,3p,twoside]{report}
\usepackage{scrextend}
\changefontsizes{13pt}
% XeLaTeX font setup (fontspec) - hỗ trợ đầy đủ tiếng Việt Unicode
\usepackage{fontspec}
\usepackage{unicode-math}
\IfFontExistsTF{Times New Roman}{%
  \setmainfont{Times New Roman}
}{%
  \setmainfont{TeX Gyre Termes}
}
% Font cho code listings (monospace) - dùng font có sẵn trên Windows
\setmonofont{Courier New}[Scale=0.9]
% Bật hyphenation cho tiếng Việt (xelatex tự xử lý qua ICU)
\usepackage{polyglossia}
\setdefaultlanguage{vietnamese}
\setotherlanguage{english}

\usepackage[top=2cm, bottom=2cm, left=3.5cm, right=2.5cm]{geometry}
\usepackage{graphicx}
\usepackage{adjustbox}  % For resizing tables and figures to fit margins

% Improve text wrapping to reduce overfull hbox warnings
\setlength{\emergencystretch}{3em}
\tolerance=1500
\usepackage{fancybox}
\usepackage{indentfirst}
\usepackage{amsthm}
\usepackage{amsmath}
\usepackage{array} % Tạo bảng array
\usepackage{enumitem} % Cho phép thay đổi kí hiệu của list
\usepackage{subfiles} % Chèn các file nhỏ, giúp chia các chapter ra nhiều file hơn
\usepackage{titlesec} % Giúp chỉnh sửa các tiêu đề, đề mục như chương, phần,..
\usepackage{titletoc}
\usepackage{chngcntr} % Dùng để thiết lập lại cách đánh số caption,..
\usepackage{pdflscape} % Đưa các bảng có kích thước đặt theo chiều ngang giấy
\usepackage{afterpage}
\usepackage[ruled,vlined]{algorithm2e}  % Hỗ trợ viết các giải thuật
\usepackage{capt-of} % Cho phép sử dụng caption lớn đối với landscape page
\usepackage{multirow} % Merge cells
\usepackage{fancyhdr} % Cho phép tùy biến header và footer
% \usepackage[natbib,backend=biber,style=ieee]{biblatex} % Giúp chèn tài liệu tham khảo
\usepackage{appendix}

\usepackage[font=small,labelfont=bf]{caption}

\usepackage{listings}
\usepackage{float}
\usepackage{subcaption}
\usepackage{xurl}

\usepackage[nonumberlist, nopostdot, nogroupskip, acronym]{glossaries}
\usepackage{glossary-superragged}
\setglossarystyle{superraggedheaderborder}
\usepackage{setspace}
\usepackage{parskip}

% package content table
\usepackage{tocbasic}

\usepackage{blindtext}


% ===================================================

\renewcommand{\bibname}{Danh_sach_tai_lieu_tham_khao} 
\usepackage[backend=bibtex,style=ieee]{biblatex}  %backend=biber is 'better'

\renewcommand\appendixname{PHỤ LỤC}
\renewcommand\appendixpagename{PHỤ LỤC}
\renewcommand\appendixtocname{PHỤ LỤC}


\addbibresource{Danh_sach_tai_lieu_tham_khao.bib} % chèn file chứa danh mục tài liệu tham khảo vào 

\include{lstlisting} % Phần này cho phép chèn code và formatting code như C, C++, Python

%\makeglossaries
\makenoidxglossaries

% Danh mục thuật ngữ và từ viết tắt

% --- Từ viết tắt tiếng Anh ---
\newglossaryentry{omr}{
    type=\acronymtype,
    name={OMR},
    description={Optical Mark Recognition - Nhận diện vạch quang học},
    first={OMR}
}
\newglossaryentry{amc}{
    type=\acronymtype,
    name={AMC},
    description={Auto Multiple Choice - Công cụ tạo đề thi trắc nghiệm tự động bằng LaTeX},
    first={AMC}
}
\newglossaryentry{api}{
    type=\acronymtype,
    name={API},
    description={Application Programming Interface - Giao diện lập trình ứng dụng},
    first={API}
}
\newglossaryentry{bloc}{
    type=\acronymtype,
    name={BLoC},
    description={Business Logic Component - Thành phần xử lý logic nghiệp vụ},
    first={BLoC}
}
\newglossaryentry{rest}{
    type=\acronymtype,
    name={REST},
    description={Representational State Transfer - Kiến trúc chuyển trạng thái đại diện},
    first={REST}
}
\newglossaryentry{jwt}{
    type=\acronymtype,
    name={JWT},
    description={JSON Web Token - Mã thông báo web JSON},
    first={JWT}
}
\newglossaryentry{rbac}{
    type=\acronymtype,
    name={RBAC},
    description={Role-Based Access Control - Kiểm soát truy cập dựa trên vai trò},
    first={RBAC}
}
\newglossaryentry{sdk}{
    type=\acronymtype,
    name={SDK},
    description={Software Development Kit - Bộ công cụ phát triển phần mềm},
    first={SDK}
}
\newglossaryentry{json}{
    type=\acronymtype,
    name={JSON},
    description={JavaScript Object Notation - Định dạng dữ liệu theo cấu trúc đối tượng JavaScript},
    first={JSON}
}
\newglossaryentry{crud}{
    type=\acronymtype,
    name={CRUD},
    description={Create, Read, Update, Delete - Tạo, đọc, cập nhật, xóa},
    first={CRUD}
}
\newglossaryentry{mvc}{
    type=\acronymtype,
    name={MVC},
    description={Model-View-Controller - Mô hình ba lớp Mô hình - Giao diện - Điều khiển},
    first={MVC}
}
\newglossaryentry{uuid}{
    type=\acronymtype,
    name={UUID},
    description={Universal Unique Identifier - Định danh duy nhất phổ quát},
    first={UUID}
}
\newglossaryentry{url}{
    type=\acronymtype,
    name={URL},
    description={Uniform Resource Locator - Địa chỉ tài nguyên thống nhất},
    first={URL}
}
\newglossaryentry{http}{
    type=\acronymtype,
    name={HTTP},
    description={HyperText Transfer Protocol - Giao thức truyền siêu văn bản},
    first={HTTP}
}
\newglossaryentry{https}{
    type=\acronymtype,
    name={HTTPS},
    description={HyperText Transfer Protocol Secure - Giao thức truyền siêu văn bản bảo mật},
    first={HTTPS}
}
\newglossaryentry{html}{
    type=\acronymtype,
    name={HTML},
    description={HyperText Markup Language - Ngôn ngữ đánh dấu siêu văn bản},
    first={HTML}
}
\newglossaryentry{css}{
    type=\acronymtype,
    name={CSS},
    description={Cascading Style Sheets - Bảng kiểu xếp chồng},
    first={CSS}
}
\newglossaryentry{sql}{
    type=\acronymtype,
    name={SQL},
    description={Structured Query Language - Ngôn ngữ truy vấn có cấu trúc},
    first={SQL}
}
\newglossaryentry{nosql}{
    type=\acronymtype,
    name={NoSQL},
    description={Not Only SQL - Không chỉ SQL},
    first={NoSQL}
}
\newglossaryentry{ui}{
    type=\acronymtype,
    name={UI},
    description={User Interface - Giao diện người dùng},
    first={UI}
}
\newglossaryentry{ux}{
    type=\acronymtype,
    name={UX},
    description={User Experience - Trải nghiệm người dùng},
    first={UX}
}
\newglossaryentry{ocr}{
    type=\acronymtype,
    name={OCR},
    description={Optical Character Recognition - Nhận diện ký tự quang học},
    first={OCR}
}
\newglossaryentry{dpi}{
    type=\acronymtype,
    name={DPI},
    description={Dots Per Inch - Số chấm trên inch},
    first={DPI}
}
\newglossaryentry{otp}{
    type=\acronymtype,
    name={OTP},
    description={One-Time Password - Mật khẩu dùng một lần},
    first={OTP}
}
\newglossaryentry{ssl}{
    type=\acronymtype,
    name={SSL},
    description={Secure Sockets Layer - Tầng ổ bảo mật giao thức},
    first={SSL}
}
\newglossaryentry{ssh}{
    type=\acronymtype,
    name={SSH},
    description={Secure Shell - Giao thức truy cập từ xa bảo mật},
    first={SSH}
}
\newglossaryentry{iot}{
    type=\acronymtype,
    name={IoT},
    description={Internet of Things - Internet vạn vật},
    first={IoT}
}
\newglossaryentry{saas}{
    type=\acronymtype,
    name={SaaS},
    description={Software as a Service - Phần mềm như dịch vụ},
    first={SaaS}
}
\newglossaryentry{docker}{
    type=\acronymtype,
    name={Docker},
    description={Docker - Nền tảng đóng gói ứng dụng dạng container},
    first={Docker}
}
\newglossaryentry{pm2}{
    type=\acronymtype,
    name={PM2},
    description={Process Manager 2 - Trình quản lý tiến trình Node.js},
    first={PM2}
}
\newglossaryentry{cnn}{
    type=\acronymtype,
    name={CNN},
    description={Convolutional Neural Network - Mạng nơ-ron tích chập},
    first={CNN}
}
\newglossaryentry{cntn}{
    type=\acronymtype,
    name={CNTT},
    description={Công nghệ thông tin},
    first={Công nghệ thông tin (CNTT)}
}
\newglossaryentry{datn}{
    type=\acronymtype,
    name={ĐATN},
    description={Đồ án tốt nghiệp},
    first={Đồ án tốt nghiệp (ĐATN)}
}


% ===================================================


\fancypagestyle{plain}{%
\fancyhf{} % clear all header and footer fields
\fancyfoot[RO,RE]{\thepage} %RO=right odd, RE=right even
\renewcommand{\headrulewidth}{0pt}
\renewcommand{\footrulewidth}{0pt}}

\setlength{\headheight}{17pt}  % Increased from 10pt to fix fancyhdr warning

\def \TITLE{ĐỒ ÁN TỐT NGHIỆP}
\def \AUTHOR{ĐINH MẠNH DŨNG}

% ===================================================
\titleformat{\chapter}[hang]{\centering\bfseries}{CHƯƠNG \thechapter.\ }{0pt}{}[]

\titleformat 
    {\chapter} % command
    [hang] % shape
    {\centering\bfseries} % format
    {CHƯƠNG \thechapter.\ } % label
    {0pt} %sep
    {} % before
    [] % after
\titlespacing*{\chapter}{0pt}{-20pt}{20pt}

\titleformat
    {\section} % command
    [hang] % shape
    {\bfseries} % format
    {\thechapter.\arabic{section}\ \ \ \ } % label
    {0pt} %sep
    {} % before
    [] % after
\titlespacing{\section}{0pt}{\parskip}{0.5\parskip}

\titleformat
    {\subsection} % command
    [hang] % shape
    {\bfseries} % format
    {\thechapter.\arabic{section}.\arabic{subsection}\ \ \ \ } % label
    {0pt} %sep
    {} % before
    [] % after
\titlespacing{\subsection}{30pt}{\parskip}{0.5\parskip}

\renewcommand\thesubsubsection{\alph{subsubsection}}
\titleformat
    {\subsubsection} % command
    [hang] % shape
    {\bfseries} % format
    {\alph{subsubsection}, \ } % label
    {0pt} %sep
    {} % before
    [] % after
\titlespacing{\subsubsection}{50pt}{\parskip}{0.5\parskip}

% \newcommand{\titlesize}{\fontsize{18pt}{23pt}\selectfont}
% \newcommand{\subtitlesize}{\fontsize{16pt}{21pt}\selectfont}
% \titleclass{\part}{top}
% \titleformat{\part}[display]
%   {\normalfont\huge\bfseries}{\centering}{20pt}{\Huge\centering}

% ===================================================
\usepackage{hyperref}
\hypersetup{pdfborder = {0 0 0}} %
\hypersetup{pdftitle={\TITLE},
	pdfauthor={\AUTHOR}}
	
\usepackage[all]{hypcap} % Cho phép tham chiếu chính xác đến hình ảnh và bảng biểu

\graphicspath{{figures/}{../figures/}} % Thư mục chứa các hình ảnh

\counterwithin{figure}{chapter} % Đánh số hình ảnh kèm theo chapter. Ví dụ: Hình 1.1, 1.2,..

\title{\bf \TITLE}
\author{\AUTHOR}

\setcounter{secnumdepth}{3} % Cho phép subsubsection trong report
% \setcounter{tocdepth}{3} % Chèn subsubsection vào bảng mục lục

\theoremstyle{definition}
\newtheorem{example}{Ví dụ}[chapter] % Định nghĩa môi trường ví dụ

\onehalfspacing
%Lùi đầu dòng
\setlength{\parindent}{15pt}

% --- Table float spacing adjustments ---
% Significantly reduce space between floats and text to eliminate large gaps
\setlength{\floatsep}{3pt plus 1pt minus 1pt}
\setlength{\textfloatsep}{5pt plus 1pt minus 1pt}
\setlength{\dblfloatsep}{3pt plus 1pt minus 1pt}
\setlength{\dbltextfloatsep}{5pt plus 1pt minus 1pt}
% In-text float separation
\setlength{\intextsep}{4pt plus 1pt minus 1pt}
% Caption spacing - reduce space above/below captions
\setlength{\abovecaptionskip}{2pt plus 1pt}
\setlength{\belowcaptionskip}{1pt}
% Reduce parskip for tighter paragraph spacing
\setlength{\parskip}{3pt}
\raggedbottom

% =========================== BODY ===============
\begin{document}
% \newgeometry{top=2cm, bottom=2cm, left=2cm, right=2cm}
\subfile{Bia} % Phần bìa
\subfile{Bia_lot} % Phần bìa
% \restoregeometry

% ===================================================
\pagenumbering{roman}
% \pagestyle{empty} % Header và footer rỗng
%\newpage
%\subfile{chapters/0_1_subject.tex}

\newpage
\pagenumbering{gobble}
\subfile{Chuong/0_2_Loi_cam_on.tex}

\newpage
\pagenumbering{gobble}
\subfile{Chuong/0_3_Tom_tat_noi_dung.tex}

% ===================================================
% \pagestyle{empty} % Header và footer rỗng
\newpage
\pagenumbering{roman} % Xóa page numbering ở cuối trang
\renewcommand*\contentsname{MỤC LỤC}

\titlecontents{chapter}
    [0.0cm]             % left margin
    {\bfseries\vspace{0.3cm}}                  % above code
    {{\bfseries{\scshape}
    CHƯƠNG \thecontentslabel.\ }}
    % numbered format
    {}         % unnumbered format
    {\titlerule*[0.3pc]{.}\contentspage}         % filler-page-format, e.g dots

    
\titlecontents{section}
    [0.0cm]             % left margin
    {\vspace{0.3cm}}                  % above code
    {\thecontentslabel \ } % numbered format
    {}         % unnumbered format
    {\titlerule*[0.3pc]{.}\contentspage}         % filler-page-format, e.g dots
    
\titlecontents{subsection}
    [1.0cm]             % left margin
    {\vspace{0.3cm}}                  % above code
    {\thecontentslabel \ } % numbered format
    {}         % unnumbered format
    {\titlerule*[0.3pc]{.}\contentspage}         % filler-page-format, e.g dots

 % Tạo mục lục tự động
\addtocontents{toc}{\protect\thispagestyle{empty}}
\phantomsection
\addcontentsline{toc}{chapter}{MỤC LỤC}
\tableofcontents
\thispagestyle{empty}
\cleardoublepage

% \pagenumbering{roman}
%Tạo danh mục hình vẽ.
\renewcommand{\listfigurename}{DANH MỤC HÌNH VẼ}
{\let\oldnumberline\numberline
\renewcommand{\numberline}{Hình~\oldnumberline}
\phantomsection
\addcontentsline{toc}{chapter}{DANH MỤC HÌNH VẼ}
\listoffigures} 
% \phantomsection\addcontentsline{toc}{section}{\numberline {} DANH MỤC HÌNH VẼ}
\newpage


 %Tạo danh mục bảng biểu.
\renewcommand{\listtablename}{DANH MỤC BẢNG BIỂU}
{\let\oldnumberline\numberline
\renewcommand{\numberline}{Bảng~\oldnumberline}
\listoftables}
% \phantomsection\addcontentsline{toc}{section}{\numberline {} DANH MỤC BẢNG BIỂU}

\glsaddall 
 \renewcommand*{\glossaryname}{Danh sách thuật ngữ}
\renewcommand*{\acronymname}{DANH MỤC THUẬT NGỮ VÀ TỪ VIẾT TẮT}
\renewcommand*{\entryname}{Thuật ngữ}
\renewcommand*{\descriptionname}{Ý nghĩa}
\printnoidxglossaries
 %\phantomsection\addcontentsline{toc}{section}{\numberline {} DANH MỤC THUẬT NGỮ VÀ TỪ VIẾT TẮT}

% \newpage
 %\subfile{chapters/0_5_Danh_muc_viet_tat.tex}

% \newpage
% \subfile{chapters/0_6_Thuat_ngu.tex}
% ===================================================


\newpage
\pagenumbering{arabic}

\pagestyle{fancy}
\fancyhf{}
\fancyhead[RE, LO]{\leftmark}
%\fancyhead[LE]{\rightmark}
\fancyfoot[RE, LO]{\thepage}

\chapter{GIỚI THIỆU ĐỀ TÀI}
\label{chapter:Introduction}
\subfile{Chuong/1_Gioi_thieu} % Phần mở đầu

\newpage
%\pagestyle{fancy} % Áp dụng header và footer
\chapter{KHẢO SÁT VÀ PHÂN TÍCH YÊU CẦU}
\label{chapter:Related_works}
\subfile{Chuong/2_Khao_sat}


\newpage
%\pagestyle{fancy} % Áp dụng header và footer
\chapter{NỀN TẢNG LÝ THUYẾT VÀ CÔNG NGHỆ SỬ DỤNG}
\label{chapter:Methodology}
\subfile{Chuong/3_Cong_nghe}

\newpage
%\pagestyle{fancy} % Áp dụng header và footer
\chapter{PHÂN TÍCH THIẾT KẾ, TRIỂN KHAI VÀ ĐÁNH GIÁ HỆ THỐNG}
\label{chapter:Experiment}
\subfile{Chuong/4_Ket_qua_thuc_nghiem}

\newpage
%\pagestyle{fancy} % Áp dụng header và footer
\chapter{CÁC GIẢI PHÁP VÀ ĐÓNG GÓP NỔI BẬT}
\label{chapter:SolutionAndContribution}
\subfile{Chuong/5_Giai_phap_dong_gop}
\newpage
%\pagestyle{fancy} % Áp dụng header và footer
\chapter{KẾT LUẬN VÀ HƯỚNG PHÁT TRIỂN} %Kết luận và hướng phát triển}
\label{chapter:conclusion}
\subfile{Chuong/6_Ket_luan}


% ===================================================
\newpage
\renewcommand\bibname{TÀI LIỆU THAM KHẢO}
\printbibliography
\phantomsection\addcontentsline{toc}{chapter}{TÀI LIỆU THAM KHẢO}

\appendixpage
\appendices
\addappheadtotoc

% \chapter*{PHỤ LỤC} %Kết luận và hướng phát triển}

%\mainmatter
\titleformat{\chapter}[hang]{\centering\bfseries}{ \thechapter.\ }{0pt}{}[]
\titlespacing*{\chapter}{0pt}{-20pt}{20pt}

\titlecontents{chapter}
    [0.0cm]             % left margin
    {\bfseries\vspace{0.3cm}}                  % above code
    {{\bfseries{\scshape} \thecontentslabel.\ }} % numbered format
    {}         % unnumbered format
    {\titlerule*[0.3pc]{.}\contentspage}         % filler-page-format, e.g dots

\chapter{ĐẶC TẢ USE CASE}
\subfile{Chuong/Phu_luc_B}

\chapter{THIẾT KẾ CHI TIẾT CÁC LỚP}
\subfile{Chuong/Phu_luc_C}

\chapter{SINH CÂU HỎI BẰNG AI}
\subfile{Chuong/Phu_luc_D}
\end{document}

